\section{Finite-algebra covers for recursive objects}
\label{sec:covers}
\status{Implemented in \code{finite_cover.py}; deterministic algebras with
constructor arities zero, one, and two; positive covers and concrete negative
certificates.}

\subsection{Replace unbounded syntax by the finite summary the goal needs}
A recursive datatype may contain arbitrarily large values while a property of
interest depends on a finite summary. Examples include counts modulo a fixed
integer, a finite-state parser's final state, a bounded flag configuration, and
equality of two folds into a finite algebra. Unrolling constructors to a fixed
depth is not a proof over all values. A finite summary with proved constructor
laws supplies a different route: compute an inductive cover of summary states.

Let $\Sigma$ be a finite ranked signature. A constructor $c$ has an arity
$a(c)$, and $\mathcal T_\Sigma$ is the set of finite ground trees built from
these constructors. A deterministic finite algebra consists of a finite set $Q$
and functions
\[
 \delta_c:Q^{a(c)}\longrightarrow Q.
\]
Every ground tree has a uniquely determined summary $\llbracket t\rrbracket$,
computed from the leaves upwards. Given $G\subseteq Q$, the decision problem is
\begin{equation}
\label{eq:tree-goal}
 \forall t\in\mathcal T_\Sigma,\quad \llbracket t\rrbracket\in G.
\end{equation}
This is a standard finite-tree-automata problem. The contribution to Forge is the
certificate form, source-level bridge, and scheduling as an alternative to
unbounded structural search. Lammich's machine-checked tree-automata development
provides a mature reference for the underlying algorithms and witness
construction, albeit in Isabelle rather than Lean~\cite{afp-tree}.

\subsection{Positive certificates are closed sets, not search logs}
\begin{definition}[Inductive cover]
A set $R\subseteq Q$ is an inductive cover if for every constructor $c$ and every
tuple $q_1,\ldots,q_{a(c)}\in R$,
\[
 \delta_c(q_1,\ldots,q_{a(c)})\in R.
\]
For a nullary constructor this requires its designated state to belong to $R$.
\end{definition}

\begin{theorem}[Finite-cover soundness]
\label{thm:finite-cover}
If $R$ is an inductive cover and $R\subseteq G$, then \eqref{eq:tree-goal} holds.
\end{theorem}
\begin{proof}
Induct on a finite constructor tree. The induction hypotheses put every child
summary in $R$. The closure check for its constructor puts the parent summary
in $R$. For nullary constructors this is exactly the required initial-state
check. Finally use $R\subseteq G$.
\end{proof}

The positive certificate can therefore be only a list of state identifiers.
The checker verifies membership, distinctness, all constructor closures, and
$R\subseteq G$. It does not need to trust how the proposer found $R$ and does
not need to recompute the least reachable set. A larger closed cover is just as
sound as the least one, provided it stays inside $G$.

\subsection{Search and a genuine counterexample}
The proposer computes the least constructor-closed set by saturation. Every
new state retains a witness tree node whose children refer to previously
constructed witness nodes. If a state outside $G$ is discovered, its witness is
returned immediately. The witness is a directed acyclic graph that denotes an
ordinary finite tree, possibly with shared subtrees.

\begin{lstlisting}
reachable := empty state-to-witness map
repeat:
    for each constructor c:
        for each tuple of already reachable child states:
            q := delta_c(tuple)
            if q is new:
                record q with a node pointing to child witnesses
                if q not in good:
                    return concrete_tree_DAG(q)
    if no new state was added:
        return inductive_cover(keys(reachable))
    if resource budget is exhausted:
        return unknown
\end{lstlisting}

The negative checker requires every child index to be strictly smaller than the
index of its parent. It evaluates each node using the input algebra and checks
that the designated root state is outside $G$. No search result tag is accepted
as a substitute for this tree. A malformed self-reference is rejected even if
its proposed summary state happens to be bad.

\begin{theorem}[Completeness for the explicit finite algebra]
Without a resource cutoff, saturation either returns an inductive cover inside
$G$ or a finite tree whose summary is outside $G$.
\end{theorem}
\begin{proof}
At most $|Q|$ new states can appear. Every stored state is realized by its
recorded finite tree. Conversely, induction on tree height shows that the
summary of every finite tree eventually enters the closure. Therefore the
fixed point is exactly the set of realizable summaries. If it contains a bad
state, the associated finite tree refutes the universal claim; otherwise it is
a positive cover.
\end{proof}

This completeness statement is for the supplied finite algebra. It is not a
claim that arbitrary recursive functions admit a sufficiently precise finite
abstraction or that finite-model reasoning decides all Lean inductive goals.

\subsection{Example: two related counts in all binary trees}
For a full binary tree let $L(t)$ be its number of leaves and $N(t)$ its number
of internal nodes. Fix $m\ge1$. Use state space
\[
 Q=(\Z/m\Z)^2,
\]
with leaf summary $(1,0)$ and node transition
\[
 (\ell_1,n_1),(\ell_2,n_2)
 \longmapsto (\ell_1+\ell_2,n_1+n_2+1).
\]
The desired good set is
\[
 G=\{(\ell,n):\ell-n\equiv1\pmod m\}.
\]
The prototype discovers this set as its reachable cover. The rule closes
because
\[
 (\ell_1+\ell_2)-(n_1+n_2+1)
 =(\ell_1-n_1)+(\ell_2-n_2)-1\equiv1.
\]
This is a theorem about \emph{all} finite binary trees, not just the trees used
to reach the $m$ states during search.

At $m=12$ the explicit product has 144 states and the binary transition table
has $144^2=20{,}736$ entries. The recorded search reaches 12 states, evaluates
145 distinct constructor/child-state tuples, and visits 235 tuples including
revisits. These are separate quantities. Building and validating the entire
explicit transition table still costs work proportional to its 20,736 entries;
145 is not an end-to-end complexity claim. A native implementation can use a
proved modular transition function instead of an explicit table, but then that
function's connection to the source fold must be checked.

The finite certificate establishes the selected congruence. It does not justify
upgrading that congruence to the integer equality $L=N+1$ merely because the
same equality motivated the choice of states.

\subsection{Example: equality of folds without assuming associativity}
Let $(A,\star)$ be an arbitrary finite magma: an operation on a finite set, with
no associativity or identity assumption. A tree has leaf labels in $A$. Compare
its fold with the fold of its mirror. The product summary is
\[
 (u,v)=\bigl(\operatorname{fold}(t),
               \operatorname{fold}(\operatorname{mirror}(t))\bigr).
\]
A leaf $a$ has summary $(a,a)$. For child summaries $(u_1,v_1),(u_2,v_2)$, the
node summary is
\[
 (u_1\star u_2,\ v_2\star v_1).
\]
The good states are the diagonal pairs. Saturation either certifies that every
reachable pair is diagonal or constructs an actual tree with unequal folds.

If $\star$ is commutative, the diagonal is closed; associativity is irrelevant.
If some $a\star b\ne b\star a$, the tree with leaves $a,b$ already provides a
counterexample. The experiment includes both classes, as well as unrelated
random algebras checked against a separately written closure oracle. This is a
small but meaningful regression suite for product construction, constructor
order, and witness reconstruction.

\subsection{The source-level bridge is a proved commuting diagram}
\label{sec:finite-bridge}
A finite table is not automatically a model of a Lean definition. For a source
datatype $D$ and summary $\alpha:D\to Q$, each constructor must have a proved
compatibility theorem
\[
 \alpha\bigl(c_D(t_1,\ldots,t_a)\bigr)
 =\delta_c\bigl(\alpha(t_1),\ldots,\alpha(t_a)\bigr).
\]
The final target must follow from $\alpha(t)\in G$. These bridge theorems are
part of preparation. With them, Theorem~\ref{thm:finite-cover} becomes a source
proof through ordinary structural induction.

For an infinite payload type, such as integer-labelled trees, a finite
classification of labels needs its own exhaustive coverage proof. A positive
cover can safely over-approximate possible labels and transitions. A bad
\emph{abstract} state does not then automatically give a concrete source
counterexample: its abstract labels or transitions might be unrealizable.
Negative export requires reconstructing actual source constructors and values
and checking the resulting tree. The current prototype uses a finite explicit
signature and exact transition tables, so it avoids this extra abstraction
layer.

A Lean implementation should begin with non-indexed datatypes and registered
summary homomorphisms. Mutual recursion, dependent indices, quotient
representations, coinductive streams, and general well-founded recursion need
additional semantics and should not be silently flattened into this signature.

\subsection{How this becomes a Forge worker}
A preparation pass recognizes a goal about a fold into \code{Fin m}, a finite
product, or another explicitly enumerated type. It searches the registered
summary laws, constructs the relevant product algebra, and sends its pure
finite description to the proposer. The returned cover is checked using
\code{Finset}/\code{Fintype} data or an array-based reflected checker. The
resulting theorem is then introduced as an ordinary local fact.

The first implementation does not need a general automata library. A worklist,
a dense or sparse state table, and the small cover checker suffice. A more
general library may become useful for nondeterministic summaries or language
operations, but importing such a library does not replace the Lean source
bridge. The reference formalization is useful as an algorithm and refinement
model, not as automatically transferable Lean authority~\cite{afp-tree}.

Use a budget in state discoveries, tuple evaluations, table size, and certificate
size. The explicit prototype implements a transition-evaluation cutoff; it is
not hardened against arbitrary hostile table sizes. State-product blow-up is the
main structural limit. Goal-directed products should include only the summaries
needed by the target, while any omission must still leave a sound abstraction.
